\begin{answers}{ch:ml}{機械学習の枠組み：データから関数を学ぶ}
\item $\calL(w)=\frac{1}{3}(14w^2-46w+38)$ で、2 ステップ後は $w\approx1.176$ です。
3 件の誤差の 2 乗 $(w-2)^2+(2w-3)^2+(3w-5)^2$ を展開して足し、3 で割ると、この式になります。
微分すると $\calL'(w)=\frac{1}{3}(28w-46)$ です。
$w=0$ では $\calL'=-15.33$ なので、$w\leftarrow0-0.05\times(-15.33)=0.767$ です。
$w=0.767$ では $\calL'=-8.18$ なので、$w\leftarrow0.767+0.05\times8.18=1.176$ です。
損失は $12.67\to3.65\to1.09$ と下がります。
谷底は $\calL'(w)=0$ となる $w=23/14\approx1.643$（損失 $0.071$）で、まだ途中です。
\item 学習率を約 $6\times10^{-5}$ 以下に下げれば発散しません（実測）。
$y$ は百トークン単位のまま、$x$ を文字数で 100 ステップ動かしました。
$\eta=6\times10^{-5}$ では損失が $0.0112$ まで下がり、$\eta=7\times10^{-5}$ では $10^{20}$ 程度に発散しました。
元の設定の境目は約 $0.40$ なので（$\eta=0.42$ で発散）、約 6,500 分の 1 です。
勾配の $w$ 成分は「誤差 $\times$ 入力」なので、入力が 100 倍になると谷が $w$ の向きにとても急になるからです。
一方、$b$ の向きの谷はゆるやかなままで、この学習率では $b$ が谷底の $0.0083$ に近づくまで約 10 万ステップかかりました。
大きさをそろえる工夫は、すべてのパラメータを同じ速さで学ばせるためにも必要です。
\item $\vp\approx(0.090,\,0.245,\,0.665)$、損失は約 $2.41$、勾配は $(-0.910,\,0.245,\,0.665)$ です。
\eqref{eq:ml-softmax}より、$e^1, e^2, e^3$ をその和 $30.19$ で割ると $\vp$ になります。
損失は正解クラスの確率の負の対数で、$-\log0.090=2.41$ です。
勾配は\eqref{eq:ml-softmax-grad}の $\vp-\vy$ で、正解のクラスの成分だけ 1 を引きます。
本文の例 $\vz=(2,1,0)$ と同じ確率の並び替えですが、正解に最も低い確率を付けたので、損失は 1.41 より大きくなります。
\item (a) $\sigmoid(z)=(1+e^{-z})^{-1}$ を合成関数の微分で微分すると、$\sigmoid'(z)=e^{-z}/(1+e^{-z})^2$ です。
これを $\frac{1}{1+e^{-z}}\times\frac{e^{-z}}{1+e^{-z}}$ と分けます。
前の因子は $\sigmoid(z)$、後の因子は $\frac{(1+e^{-z})-1}{1+e^{-z}}=1-\sigmoid(z)$ です。
(b) softmax の式で $K=2$ とすると $p_1=e^{z_1}/(e^{z_1}+e^{z_2})$ です。
分母と分子を $e^{z_1}$ で割ると $p_1=1/(1+e^{-(z_1-z_2)})=\sigmoid(z_1-z_2)$ になります。
2 クラスの softmax 回帰は、ロジットの差を一つのロジットとみなしたロジスティック回帰と同じです。
\item 誤差の絶対値の和（平均）を最小にすることと同じです。
1 件分の対数は $-\log(2s)-\abs{y-f_\theta(x)}/s$ で、
$N$ 件分足して符号を変えると、$N\log(2s)+\frac{1}{s}\sum_i\abs{y^{(i)}-f_\theta(x^{(i)})}$ になります。
第 1 項はパラメータを含まず、$\frac{1}{s}$ は正の定数なので、\eqref{eq:ml-gauss-nll}の 2 乗が絶対値に置き換わっただけです。
絶対値は大きな外れを 2 乗ほど重く数えないので、ずれにときどき大きな外れがあると考えるなら、こちらが合います。
\item 最小になる次数は $M=3$ から $M=6$ に変わりますが、$M=3$〜$6$ の差はごくわずかです（実測）。
訓練データの件数を 20 にして、同じスクリプトを実行しました。
検証誤差は $M=3$ で $0.281$、$M=5$ で $0.279$、$M=6$ で $0.279$（最小）、$M=9$ で $0.291$ でした。
差は $0.003$ 程度なので、「3〜6 次がほぼ横並びになった」と読むのが正確です（乱数の種を変えると最小の次数は 3〜6 の間で入れ替わりました）。
はっきり変わったのは高い次数で、$M=9$ の検証誤差は $0.459$ から $0.291$ に下がりました。
データが増えると、複雑なモデルでも過学習しにくくなります。
\end{answers}
